\chapter{松弛内存一致性}

\begin{epigraph}
\textbf{核心命题}　现代 RISC（RISC-V RVWMO、IBM POWER、ARMv8）普遍采用松弛一致性模型，允许编译器与硬件大幅重排内存操作以隐藏延迟。本章讨论 XC 示例模型、释放一致性（RC）、数据竞争自由（DRF）契约，以及商用松弛模型的差异。
\end{epigraph}



\section{动机}

\subsection{重排序内存操作的机会}

SC 与 TSO 都限制了内存操作的重排。但现代处理器有大量机制可以受益于更激进的重排：

\begin{itemize}
\item \textbf{非阻塞缓存}（non-blocking cache）：miss 时不阻塞后续 hit；
\item \textbf{乱序执行}（out-of-order execution）：load/store 可以乱序执行；
\item \textbf{写合并}（write coalescing）：多个 store 合并为一次写；
\item \textbf{预取}（prefetch）：提前取数据。
\end{itemize}


这些机制在 SC/TSO 下被部分限制，但\textbf{如果放松更多重排}，性能可以进一步提升。

\subsection{利用重排序的机会}

松弛模型的核心思想：\textbf{默认允许所有重排，仅在程序员显式要求同步时才施加顺序约束}。这通过 FENCE 指令或 acquire/release 语义实现。代价是程序员负担加重——必须正确插入 FENCE，否则会出现难以调试的并发 bug。

\section{一个松弛一致性模型（XC）示例}

Sorin/Hill/Wood 教材定义了一个名为 \textbf{XC}（Relaxed, Coherent）的示例模型，它具备松弛模型的所有关键特征但又不至于过于复杂。

\subsection{XC 模型的基本概念}

XC 允许以下重排（不同地址之间）：

\begin{itemize}
\item load$\rightarrow$store
\item load$\rightarrow$load
\item store$\rightarrow$store
\item store$\rightarrow$load（与 TSO 相同）
\end{itemize}


但对\textbf{同一地址}的操作仍保持连贯性约束（load 读最近 store）。XC 只通过 FENCE 恢复顺序。

\subsection{使用 XC 的 FENCE 的示例}

考虑发布模式：

\begin{verbatim}
P1:                        P2:
data = 42;                 while (flag == 0);
FENCE;                     FENCE;
flag = 1;                  print(data);
\end{verbatim}


FENCE 保证 P1 的 \texttt{data = 42} 在 \texttt{flag = 1} 之前对其他核可见，P2 的 load(flag) 在 load(data) 之前。没有 FENCE，XC 允许这些重排，发布模式会失效。

\subsection{形式化 XC}

XC 的形式化用"保持程序顺序"（preserved program order, PPO）关系刻画——只有被 FENCE 标记的顺序被保持，其余可重排。

\subsection{展示 XC 正确运行的示例}

XC 的正确性体现为：在正确使用 FENCE 的程序中，XC 的行为等价于 SC。这与下一节的 DRF 契约紧密相关。

\section{实现 XC}

XC 的实现比 TSO 更激进：

\begin{itemize}
\item store buffer 可以满负荷工作；
\item 不同地址的 load/store 可以乱序；
\item FENCE 通过 drain buffer + 等待 ack 实现。
\end{itemize}


\subsection{在 XC 中使用原子指令}

原子指令（如 LR/SC、compare-and-swap）在 XC 中自带 FENCE 语义——它们既是原子的，也施加 acquire/release 顺序。

\subsection{在 XC 中的 FENCE}

FENCE 在 XC 中是全屏障，drain 所有挂起的 load/store 并等待 ack。

\subsection{一个警告}

松弛模型的威力也是它的危险——\textbf{漏掉一个 FENCE 就会导致难以复现的并发 bug}。这正是 DRF 契约与高级语言内存模型（C++/Java）的价值所在——它们把 FENCE 的正确使用抽象到语言层。

\section{无数据竞争程序的连续一致性}

\textbf{数据竞争自由（DRF, Data-Race-Free）契约}是松弛模型的理论基石，由 Adve \& Hill 在 1990 年 ISCA 论文《Weak Ordering – A New Definition》中提出：

\begin{epigraph}
如果一个程序是无数据竞争的（即所有冲突访问都被同步操作保护），那么它在松弛模型下的行为等价于 SC。
\end{epigraph}


这就是著名的 \textbf{SC for DRF} 契约。它的意义在于：程序员只需保证程序无数据竞争（用锁、原子操作正确同步），就能享受 SC 的直观语义，同时获得松弛模型的性能。

这个契约是 C++11、Java 内存模型的理论基础——高级语言通过 SC-DRF 向程序员承诺"正确同步的程序行为等价于 SC"。

\section{一些松弛模型概念}

\subsection{释放一致性}

\textbf{释放一致性}（Release Consistency, RC）由 Gharachorloo 等人在 1990 年 ISCA 论文《Memory Consistency and Event Ordering in Scalable Shared-Memory Multiprocessors》中提出。RC 区分两类同步操作：

\begin{itemize}
\item \textbf{acquire}（获取，如 lock）：acquire 之后的普通访问不能重排到 acquire 之前；
\item \textbf{release}（释放，如 unlock）：release 之前的普通访问不能重排到 release 之后；
\item acquire 与 release 之间的普通访问可以任意重排。
\end{itemize}


RC 比 XC 更精细——它把同步语义编码到 acquire/release 标注里，减少了对显式 FENCE 的依赖。RC 影响了 Stanford DASH 多处理器，被引超 1800 次。

\subsection{因果性和写原子性}

松弛模型需要维护两个关键不变量：

\begin{itemize}
\item \textbf{写原子性}（write atomicity / multicopy atomicity）：一个 store 被所有核同时观察到。否则不同观察者可能看到不一致顺序（IRIW 反例）。
\item \textbf{因果性}（causality）：禁止"自因果"循环——一个写不能由它自身推断出来（Java 内存模型特别关注这点）。
\end{itemize}


\section{松弛内存模型案例研究}

\subsection{RVWMO}

\textbf{RISC-V Weak Memory Ordering（RVWMO）}是 RISC-V 的默认内存模型，载于 RISC-V Unprivileged ISA Manual。RVWMO 的核心是\textbf{保持程序顺序（PPO, Preserved Program Order）}——共 12 条规则，分四组：

\begin{itemize}
\item \textbf{重叠地址规则}（3 条）：同地址操作的顺序；
\item \textbf{显式同步规则}（4 条）：FENCE、acquire（aq 位）、release（rl 位）、RCsc 配对；
\item \textbf{句法依赖规则}（3 条）：地址/数据/控制依赖；
\item \textbf{流水线依赖规则}（2 条）：经中间 store 的传递依赖。
\end{itemize}


RVWMO 的 aq/rl 位让 AMO/LR/SC 指令携带 acquire/release 语义，无需额外 FENCE。

\subsection{IBM POWER}

IBM POWER 是最松的商用模型之一（在 ARMv8 之前），允许 store$\rightarrow$load、store$\rightarrow$store、load$\rightarrow$load、load$\rightarrow$store 全部重排。三类屏障：

\begin{itemize}
\item \texttt{sync}（重权，全屏障）；
\item \texttt{lwsync}（轻权，不排 store$\rightarrow$load）；
\item \texttt{isync}（指令同步，控制依赖）。
\end{itemize}


POWER 的形式化由 Sarkar 等人在 PLDI 2011 论文《Understanding POWER Multiprocessors》中给出（被引超 380 次）。

\section{进一步阅读和商用松弛内存模型}

\subsection{学术文献}

\begin{itemize}
\item Adve \& Gharachorloo 1996 综述（经典）；
\item Alglave et al. PLDI 2014《Herding Cats》（公理化统一框架，被引超 530）；
\item Maranget/Sarkar/Sewell《A Tutorial Introduction to the ARM and POWER Relaxed Memory Models》。
\end{itemize}


\subsection{商用模型}

ARMv8 从 ARMv7 的非多副本原子模型修订为 multicopy-atomic（Pulte 等 POPL 2018，被引超 240）。RISC-V RVWMO 是 multicopy-atomic 的。IBM POWER 是非 multicopy-atomic 的（最松）。

\section{比较内存模型}

\subsection{松弛内存模型彼此之间以及与 TSO 和 SC 的关系}

严格程度排序（从严格到松）：SC > TSO > ARMv8/RVWMO > IBM POWER。

SC 禁止所有重排；TSO 仅允许 store$\rightarrow$load；ARMv8/RVWMO 允许大部分重排但保留 acquire/release；POWER 几乎允许所有重排。

\subsection{松弛模型有多好}

松弛模型在性能上明显优于 SC/TSO——尤其在多核可扩展性与隐藏延迟方面。代价是程序员负担与验证复杂度。SC-DRF 契约在很大程度上缓解了程序员负担。

\section{高级语言模型}

\subsection{C++11 内存模型}

C++11（ISO/IEC 14882:2011）引入了 \texttt{std::memory\_order}，由 Boehm \& Adve 在 PLDI 2008《Foundations of the C++ Concurrency Memory Model》奠基（被引超 650）。核心是 SC-DRF 契约。\texttt{memory\_order} 枚举：

\begin{itemize}
\item \texttt{relaxed}：无序；
\item \texttt{consume}：数据依赖（C++17 起建议不用）；
\item \texttt{acquire}：后续读写不可重排到 acquire 之前；
\item \texttt{release}：此前读写不可重排到 release 之后；
\item \texttt{acq\_rel}：acquire + release；
\item \texttt{seq\_cst}：额外保证单一全局总序（最强）。
\end{itemize}


\subsection{Java 内存模型}

Java 内存模型（JLS §17.4）由 Manson、Pugh、Adve 在 POPL 2005《The Java Memory Model》中提出。它用 happens-before + 因果性要求（causality requirements）刻画合法执行，纠正了原 JVM 模型的安全漏洞。

\section{小结}

松弛模型允许编译器与硬件大幅重排内存操作，仅在显式同步时施加顺序约束。XC 示例模型、释放一致性（RC）是代表性框架。DRF 契约（SC for DRF）是松弛模型的理论基石——无数据竞争的程序在松弛模型下行为等价于 SC。RISC-V RVWMO（13 条 PPO 规则 + aq/rl 位）与 IBM POWER（最松）是两端。C++11/Java 内存模型通过 SC-DRF 向程序员承诺直观语义。松弛模型在性能上明显优于 SC/TSO，但需要程序员正确使用 FENCE/acquire/release。

\section*{参考文献}
\begin{itemize}
\item Adve S. V., Hill M. D. Weak ordering – A new definition. ISCA 1990, pp. 2–14.
\item Adve S. V., Hill M. D. A unified formalization of four shared-memory models. \emph{IEEE TPDS}, 1993, 4(6): 613–624. http://rsim.cs.uiuc.edu/Pubs/ps2pdf/tpds93.pdf
\item Gharachorloo K., Lenoski D., Laudon J., et al. Memory consistency and event ordering in scalable shared-memory multiprocessors. ISCA 1990, pp. 15–26. https://doi.org/10.1145/325096.325102
\item Boehm H.-J., Adve S. V. Foundations of the C++ concurrency memory model. PLDI 2008, pp. 68–78. https://doi.org/10.1145/1375581.1375591
\item Manson J., Pugh W., Adve S. V. The Java memory model. POPL 2005, pp. 378–391. https://rsim.cs.uiuc.edu/Pubs/popl05.pdf
\item Sarkar S., Sewell P., et al. Understanding POWER multiprocessors. PLDI 2011. https://www.cl.cam.ac.uk/\textasciitilde{}pes20/ppc-supplemental/pldi105-sarkar.pdf
\item Pulte C., Flur S., Deacon W., et al. Simplifying ARM concurrency: Multicopy-atomic axiomatic and operational models for ARMv8. POPL 2018. https://doi.org/10.1145/3158107
\item Alglave J., Maranget L., Tautschnig M. Herding cats. PLDI 2014. http://diy.inria.fr/herd/herding-cats-color.pdf
\item RISC-V. \emph{Unprivileged ISA Manual, Chapter 17: RVWMO}. https://docs.riscv.org/reference/isa/v20260120/unpriv/rvwmo.html
\item Sorin D. J., Hill M. D., Wood D. A. \emph{A Primer on Memory Consistency and Cache Coherence} (2nd ed.). 2020, Chapter 5. https://pages.cs.wisc.edu/\textasciitilde{}markhill/papers/primer2020\_2nd\_edition.pdf
\item Maranget L., Sarkar S., Sewell P. A tutorial introduction to the ARM and POWER relaxed memory models. https://www.cl.cam.ac.uk/\textasciitilde{}pes20/ppc-supplemental/test7.pdf
\item cppreference. \emph{std::memory\_order}. https://en.cppreference.com/cpp/atomic/memory\_order
\item Pugh W. \emph{The Java Memory Model FAQ}. https://www.cs.umd.edu/\textasciitilde{}pugh/java/memoryModel/jsr-133-faq.html
\end{itemize}
